This paper develops an ontology, a phenomenology, and an ethics of culture as it emerges in intimate relation. Its guiding claim is that culture is misconceived when taken as a finished and transmissible object and is rightly conceived as a generation that does not close: the dialectical exteriorization, across the three registers of experience and under material constraint, of the interwoven generative history of relational subjects. The paper works first at the level of culture in general, passing through the major theories of cultural formation---transmission and collective representation, symbol and interpretation, structure and reproduction, culture as the encoding of power, the Chinese \emph{wen--hua} tradition, and the dialogic---and locating within them a position they leave open. That position is stated as a variant of the Lacanian theory of the three registers, here called the \emph{three-register dialectic of the relational Real}: the Real is granted a hitherto neglected relational form; culture is shown to span all three registers, none dispensable; and the registers are set into a materially driven, dialectical motion, so that culture has no final truth, is never static, and develops through the historical sublation of its own forms.

On this ground the paper turns to its proper subject, the emergence of culture in the intimate dyad, which it takes as the minimal and most exposed unit of cultural emergence. Five mechanisms are shown to belong to the dyad alone: the vacancy of any external authority to fix meaning, so that the two must serve as each other's symbolic authority; the mutual interweaving of two structures of desire; a density of shared living that lets culture sediment at a visible speed; the heightened, intimate configuration of the three registers; and a fragility by which a whole shared culture may collapse. A middle-scale case, the family ethos (\emph{jiafeng}), is examined to draw out the epistemology of a dialectical view of culture and the study of power within it. The paper closes with the political economy and ethics of intimate culture: how such culture is produced under material conditions, how capital, having exhausted older frontiers, reaches into the bond itself to capture the value it generates, and on what the just generation of culture depends---namely, that the generation of value, which carries power, must in the same movement generate the counter-power that restrains power and prevents its entrenchment.

Beyond its revision of Lacan, the paper advances five theses of its own: the \emph{non-transferability thesis}, that a relational culture's resistance to transfer rises with the depth of its anchoring in interwoven drive; the \emph{vacancy thesis}, that the externalization of symbolic authority sets a system's endogeneity and its stability in inverse tension; \emph{capture-as-decoherence}, that capital captures a relational culture not by seizing it but by preserving its symbolic form while drawing off its coherence in the Real; the \emph{counter-power criterion}, that a cultural form persists through changed conditions only insofar as it endogenously generates a counter-power against its own power; and the \emph{real-anchoring thesis}, that a relational culture's resistance to power, capital, and time arises from its anchoring in the register of the Real.
